% ============================================================
%  DEAMs 24h — Apresentação executiva (6 slides)
%  Estilo editorial: capa NAVY + conteúdo CREME, serif + sans,
%  acentos dourados. Inspirado no template "VOL · TRADING".
%  COMPILAR COM XeLaTeX (rodar 2x):  xelatex deam24h_slides.tex
% ============================================================
\documentclass[aspectratio=169,11pt]{beamer}

% ── Fontes do sistema (Windows) via fontspec/XeLaTeX ────────
\usepackage{fontspec}
\setmainfont{Georgia}      % SERIF  -> títulos/headings (\rmfamily)
\setsansfont{Calibri}      % SANS   -> corpo/labels  (\sffamily)
\renewcommand{\familydefault}{\sfdefault}   % corpo = Calibri
\usepackage[brazil]{babel}
\usepackage{amsmath, amssymb}

% ── Gráficos, tabelas, caixas, diagramas ────────────────────
\usepackage{graphicx}
\graphicspath{{../figuras_causal/}}
\usepackage{booktabs}
\usepackage{array}
\usepackage{colortbl}
\usepackage{tikz}
\usetikzlibrary{shadows}
\usepackage[most]{tcolorbox}

% ── Paleta (extraída do template de referência) ─────────────
\definecolor{navy}{HTML}{0E1B33}    % fundo da capa
\definecolor{navy2}{HTML}{16294C}   % títulos/headings (serif)
\definecolor{cream}{HTML}{F4F1E8}   % fundo do conteúdo
\definecolor{creamtx}{HTML}{ECE7D8} % texto claro sobre navy
\definecolor{gold}{HTML}{C19A3E}    % acento dourado
\definecolor{ink}{HTML}{2C2C2C}     % corpo de texto
\definecolor{slate}{HTML}{5C6B82}   % itálico/subtítulo
\definecolor{card}{HTML}{FBFAF5}    % fundo de card
\definecolor{hair}{HTML}{D8D2C2}    % filetes
\definecolor{ared}{HTML}{A8443C}    % alerta (vermelho terroso)
\definecolor{agreen}{HTML}{4F7A4E}  % ok/solução (verde sóbrio)

% ── Tema base ───────────────────────────────────────────────
\usetheme{default}
\setbeamertemplate{navigation symbols}{}
\setbeamercolor{background canvas}{bg=cream}
\setbeamercolor{normal text}{fg=ink, bg=cream}
\setbeamercolor{itemize item}{fg=gold}
\setbeamerfont{frametitle}{family=\rmfamily}
\setbeamertemplate{frametitle}{}            % não usamos o frametitle nativo
% marcador de lista = traço dourado (igual à referência)
\setbeamertemplate{itemize item}{\color{gold}\rule[0.28ex]{0.85em}{2.2pt}}
\setbeamertemplate{itemize subitem}{\color{gold}\textbullet}

% ── Rodapé editorial (ponto dourado + label + nº de página) ─
\setbeamertemplate{footline}{%
  \hbox to\paperwidth{\hskip0.9cm
    {\color{gold}\footnotesize\textbullet}\hskip0.4em
    {\sffamily\scriptsize\bfseries\color{slate}%
      \MakeUppercase{Avaliação de impacto\, $\cdot$\, DEAMs 24h no Brasil}}%
    \hfill
    {\rmfamily\itshape\footnotesize\color{slate}%
      \ifnum\value{framenumber}<10 0\fi\insertframenumber\,/\,06}%
    \hskip0.9cm}%
  \vspace{3pt}%
}

% ── Cabeçalho do slide (chamar no topo de cada frame) ───────
% \slidehead{KICKER}{Título}{Subtítulo}
\newcommand{\slidehead}[3]{%
  {\sffamily\scriptsize\bfseries\color{gold}\textbullet\hskip0.45em
    \color{navy2}\MakeUppercase{#1}}%
  \hfill
  {\rmfamily\itshape\footnotesize\color{slate}DEAM $\cdot$ BR}\\[2pt]
  {\color{hair}\rule{\textwidth}{0.6pt}}\\[6pt]
  \begin{tcolorbox}[enhanced, colback=cream, colframe=cream, boxrule=0pt,
       borderline west={3pt}{0pt}{gold}, left=12pt, right=0pt,
       top=0pt, bottom=0pt, boxsep=0pt, before skip=0pt, after skip=0pt,
       width=\textwidth]
    {\rmfamily\bfseries\color{navy2}\Large #2}\\[3pt]
    {\rmfamily\itshape\color{slate}\small #3}
  \end{tcolorbox}%
  \vspace{9pt}
}

% altura máxima padrão das figuras (cabe na área útil do slide)
\newcommand{\figmax}[1]{\includegraphics[width=\linewidth,height=4.3cm,keepaspectratio]{#1}}

% ── Eyebrow dourado (mini-rótulo de seção) ──────────────────
\newcommand{\eyebrow}[1]{%
  {\sffamily\scriptsize\bfseries\color{gold}\MakeUppercase{#1}}\par\vspace{5pt}}

% ── Item de "leitura" (ponto dourado + termo + descrição) ───
\newcommand{\readitem}[2]{%
  \noindent{\color{gold}\textbullet}\hskip0.55em
  {\sffamily\bfseries\color{navy2}#1}%
  {\sffamily\color{ink}\ \ #2}\par\vspace{5pt}}

% ── Card numerado (estilo referência) ───────────────────────
\newtcolorbox{numcard}[2]{enhanced, colback=card, colframe=hair,
  boxrule=0.5pt, arc=2pt, drop fuzzy shadow={black!18},
  borderline north={2pt}{0pt}{gold},
  left=10pt, right=10pt, top=8pt, bottom=8pt,
  fonttitle=\rmfamily,
  title={{\rmfamily\itshape\large\color{gold}#1}\\[2pt]
         {\rmfamily\bfseries\color{navy2}\large #2}},
  coltitle=navy2, attach title to upper=\par\vspace{6pt}}

% ── Callout colorido (veredicto/intuição) ───────────────────
\newtcolorbox{callout}[2][gold]{enhanced, colback=card, colframe=#1,
  boxrule=0pt, leftrule=3pt, arc=2pt,
  left=9pt, right=9pt, top=4pt, bottom=4pt,
  before skip=4pt, after skip=2pt,
  fonttitle=\sffamily\bfseries\footnotesize, coltitle=#1,
  title={\MakeUppercase{#2}}}

% ── Card branco para gráficos ───────────────────────────────
\newtcolorbox{chartcard}{enhanced, colback=white, colframe=hair,
  boxrule=0.5pt, arc=2pt, drop fuzzy shadow={black!15},
  left=4pt, right=4pt, top=4pt, bottom=4pt}

% ── Passo da cadeia causal ──────────────────────────────────
\newtcolorbox{chainstep}[1][navy2]{enhanced, colback=card, colframe=hair,
  boxrule=0.5pt, arc=2pt, borderline west={2.5pt}{0pt}{#1},
  left=8pt, right=8pt, top=2pt, bottom=2pt,
  before skip=0pt, after skip=0pt, halign=center}
\newcommand{\chainarrow}{\vspace{1pt}\centerline{\color{gold}$\downarrow$}\vspace{1pt}}

% ============================================================
\begin{document}

% ╔══════════════════════════════════════════════════════════╗
% ║  SLIDE 0 — CAPA (fundo navy)                              ║
% ╚══════════════════════════════════════════════════════════╝
{%
\setbeamercolor{background canvas}{bg=navy}
\begin{frame}[plain]
  % filetes verticais dourados na borda direita
  \begin{tikzpicture}[overlay, remember picture]
    \draw[gold, line width=0.5pt]
      ([xshift=-1.7cm]current page.north east) -- ([xshift=-1.7cm]current page.south east);
    \draw[gold, line width=0.5pt]
      ([xshift=-1.3cm]current page.north east) -- ([xshift=-1.3cm]current page.south east);
  \end{tikzpicture}

  {\color{gold}\footnotesize\textbullet}\hskip0.4em
  {\rmfamily\itshape\footnotesize\color{creamtx}DEAM $\cdot$ BR}

  \vfill
  {\sffamily\small\bfseries\color{gold}\MakeUppercase{Avaliação de Impacto Causal}}

  \vspace{12pt}
  {\rmfamily\bfseries\color{creamtx}\fontsize{30}{36}\selectfont
   Conversão de DEAMs para\\ Plantão 24h no Brasil}

  \vspace{14pt}
  {\color{gold}\rule{2.2cm}{1.6pt}}

  \vspace{14pt}
  {\rmfamily\itshape\color{creamtx}\large
   Do acesso à proteção --- uma abordagem via Callaway \& Sant'Anna (2021)}

  \vspace{10pt}
  {\sffamily\color{creamtx}\normalsize
   Felipe Silva \;{\color{gold}$\cdot$}\; Rafael Issa de Lima}
  \vfill

  {\sffamily\scriptsize\bfseries\color{gold}%
   \MakeUppercase{Apresentação técnica\, $\cdot$\, DEAM-BR\, $\cdot$\, Painel 2009--2019}}
  \vspace{4pt}
\end{frame}
}

% ╔══════════════════════════════════════════════════════════╗
% ║  SLIDE 1 — O PROBLEMA E A CADEIA CAUSAL                   ║
% ╚══════════════════════════════════════════════════════════╝
\begin{frame}[t]
  \slidehead{Contexto $\cdot$ 01}{O Problema e a Cadeia Causal}%
            {Por que ampliar o horário de atendimento das delegacias?}

  \begin{columns}[T]
    \begin{column}{0.50\textwidth}
      \eyebrow{O problema}
      {\sffamily\color{ink}\small DEAMs operam, em regra, em
       \textbf{horário comercial} (9h--18h). Mas a violência não:}

      \vspace{14pt}
      \begin{flushleft}
        {\rmfamily\bfseries\color{ared}\fontsize{58}{60}\selectfont 57\%}\\[4pt]
        {\sffamily\color{slate}\small das notificações ocorrem \textbf{fora}
         do horário comercial (08h--17h)}
      \end{flushleft}
    \end{column}

    \begin{column}{0.50\textwidth}
      \eyebrow{Cadeia causal (cifra oculta)}
      \begin{chainstep}[gold]
        {\rmfamily\bfseries\color{navy2}Plantão 24h}
      \end{chainstep}
      \chainarrow
      \begin{chainstep}[agreen]
        {\rmfamily\bfseries\color{navy2}$\uparrow$ Acesso}\\
        {\sffamily\scriptsize\color{slate}mais notificações noturnas}
      \end{chainstep}
      \chainarrow
      \begin{chainstep}[agreen]
        {\rmfamily\bfseries\color{navy2}Proteção}\\
        {\sffamily\scriptsize\color{slate}medidas protetivas de urgência}
      \end{chainstep}
      \chainarrow
      \begin{chainstep}[ared]
        {\rmfamily\bfseries\color{navy2}$\downarrow$ Letalidade}\\
        {\sffamily\scriptsize\color{slate}feminicídios}
      \end{chainstep}
    \end{column}
  \end{columns}
\end{frame}

% ╔══════════════════════════════════════════════════════════╗
% ║  SLIDE 2 — METODOLOGIA E DESAFIO DOS DADOS                ║
% ╚══════════════════════════════════════════════════════════╝
\begin{frame}[t]
  \slidehead{Metodologia $\cdot$ 02}{Metodologia e o Desafio dos Dados}%
            {Identificação sob adoção escalonada --- desfechos em taxa /100k hab.}

  \begin{columns}[T]
    \begin{column}{0.50\textwidth}
      \begin{callout}[ared]{O problema do TWFE}
        {\sffamily\footnotesize Goodman-Bacon (2021): com \textbf{adoção
        escalonada} (10 coortes), o estimador de efeitos fixos gera
        \textbf{pesos negativos} sob efeitos heterogêneos.}
      \end{callout}
      \vspace{6pt}
      \begin{callout}[agreen]{A solução $\cdot$ CS-DiD (2021)}
        {\sffamily\footnotesize Efeitos por coorte-tempo $ATT(g,t)$: os
        \textbf{38 tratados} comparados \textbf{apenas} com os
        \textbf{247 nunca-tratados} (grupo de controle limpo).}
      \end{callout}
    \end{column}

    \begin{column}{0.50\textwidth}
      \eyebrow{38 tratados em 10 coortes (2014 = maior)}
      \begin{chartcard}
        \figmax{00b_coortes.png}
      \end{chartcard}
    \end{column}
  \end{columns}
\end{frame}

% ╔══════════════════════════════════════════════════════════╗
% ║  SLIDE 3 — RESULTADOS I: ACESSO                           ║
% ╚══════════════════════════════════════════════════════════╝
\begin{frame}[t]
  \slidehead{Resultados $\cdot$ 03}{Acesso e o Viés de Tendência}%
            {Desfecho: taxa de notificações /100k (SINAN)}

  \begin{columns}[T]
    \begin{column}{0.58\textwidth}
      \begin{chartcard}
        \figmax{notificacoes_02_event_study.png}
      \end{chartcard}
    \end{column}

    \begin{column}{0.42\textwidth}
      \eyebrow{Leitura científica}
      \readitem{ATT $= -20{,}51$}{($p=0{,}001$)}
      \readitem{Pré-tendências violadas}{$p_{\text{Wald}}=0{,}004$}
      \readitem{Pontos em $t<0$}{já divergem de zero}

      \vspace{4pt}
      \begin{callout}[ared]{Veredicto}
        {\sffamily\small O efeito capta uma \textbf{trajetória prévia},
        não a política. \textbf{Não interpretável como causal.}}
      \end{callout}
    \end{column}
  \end{columns}
\end{frame}

% ╔══════════════════════════════════════════════════════════╗
% ║  SLIDE 4 — RESULTADOS II: LETALIDADE                      ║
% ╚══════════════════════════════════════════════════════════╝
\begin{frame}[t]
  \slidehead{Resultados $\cdot$ 04}{Letalidade e Adoção Reativa}%
            {Desfecho: taxa de feminicídios /100k (SIM)}

  \begin{columns}[T]
    \begin{column}{0.58\textwidth}
      \begin{chartcard}
        \figmax{feminicidios_02_event_study.png}
      \end{chartcard}
    \end{column}

    \begin{column}{0.42\textwidth}
      \eyebrow{Leitura científica}
      \readitem{Pré-tendências OK}{$p_{\text{Wald}}=0{,}165$}
      \readitem{ATT $= +0{,}438$}{($p=0{,}082$) --- sinal invertido}

      \vspace{4pt}
      \begin{callout}[gold]{Intuição $\cdot$ adoção reativa}
        {\sffamily\small Gestores abrem DEAMs 24h \textbf{onde a letalidade
        já está em viés de alta} --- causalidade reversa
        (\textit{endogeneidade}).}
      \end{callout}
    \end{column}
  \end{columns}
\end{frame}

% ╔══════════════════════════════════════════════════════════╗
% ║  SLIDE 5 — MECANISMO E ROADMAP                            ║
% ╚══════════════════════════════════════════════════════════╝
\begin{frame}[t]
  \slidehead{Mecanismo \& agenda $\cdot$ 05}{Mecanismo e o Caminho à Frente}%
            {Evidência de sazonalidade e próximos passos metodológicos}

  \begin{columns}[T]
    \begin{column}{0.50\textwidth}
      \eyebrow{A demanda noturna é constante}
      {\sffamily\small\color{ink}
      \setlength{\arrayrulewidth}{0.4pt}\arrayrulecolor{hair}
      \renewcommand{\arraystretch}{1.05}
      \begin{tabular}{@{}p{0.62\linewidth}>{\raggedleft\arraybackslash}p{0.26\linewidth}@{}}
        \toprule
        {\bfseries\color{navy2}Grupo / Período} & {\bfseries\color{navy2}\% fora 08--17h}\\
        \midrule
        Tratados --- antes do 24h & \color{ared}57,1\%\\
        Tratados --- após o 24h   & \color{ared}55,9\%\\
        Controles (comercial)     & \color{ared}57,6\%\\
        \bottomrule
      \end{tabular}}

      \vspace{6pt}
      {\rmfamily\itshape\small\color{slate} $\sim$\textbf{\color{navy2}56\%}
       fora do horário comercial \textbf{\color{navy2}antes e depois} do 24h.}
    \end{column}

    \begin{column}{0.50\textwidth}
      \eyebrow{Roadmap $\cdot$ próximos passos}
      \begin{itemize}\sffamily\small
        \item \textbf{DR-DiD} --- estimador duplamente robusto.
        \item \textbf{Covariáveis socioeconômicas} (IDH, renda, urbanização)
              $\Rightarrow$ tendências paralelas \textit{condicionais}.
        \item \textbf{Controle de antecipação} e grupo de
              não-ainda-tratados.
      \end{itemize}

      \vspace{8pt}
      \begin{callout}[navy2]{Mensagem central}
        {\sffamily\small Sem ajustar para \textbf{adoção reativa} e
        \textbf{tendências prévias}, avaliações de plantão 24h tendem a ser
        \textbf{pessimistas} por construção.}
      \end{callout}
    \end{column}
  \end{columns}
\end{frame}

\end{document}
